\section{Horospheres}

The claim that our flat three-dimensional space is the cusp of the curved four-dimensional crystal needs one careful qualification, because it is easy to hear wrong. It does not say that far enough out the bulk turns flat. The bulk is hyperbolic everywhere and forever, at every distance, with no edge where the curvature gives out. What is flat is the geometry on a particular surface through that curved space, a horosphere, and its flatness is exact.

A horosphere is a sphere whose center has been pushed all the way to infinity, an infinite-radius sphere, so its surface no longer curves back on itself but lies like an endless sheet inside the bulk. In two dimensions it is the horocycle resting tangent to the rim of the Poincar\'e disk. A small patch of the round Earth only looks flat because the globe is large, but push the radius to infinity and the surface is not nearly flat, it is exactly flat. So is a horosphere. Its intrinsic geometry is Euclidean not in a limit but at every finite distance, a theorem and not a hope. In the natural coordinates the line element is $ds^2 = dr^2 + e^{2r}\,dx^2$, with $r$ the outward radial direction and $x$ the directions along the surface. Fix $r$ to stand on one horosphere and the induced measure is $e^{2r}$ times the flat Euclidean one, a plain change of scale that leaves it flat. The radial direction crossing the sheet stays hyperbolic, but the sheet itself is Euclidean three-space, exactly, however close in you take it. Flatness, then, is not the thing that needs infinity.

What needs infinity is only the perfect, endless cubic crystal of cells, the complete $\{4,3,4\}$ honeycomb wrapped around the ideal point. And that is an idealization in exactly the sense solid-state physics already lives with, for no real crystal is infinite either, yet the infinite-lattice picture is used because a large finite chunk is locally identical to it and its bulk properties settle long before any edge. The cusp is the same. The finite mesh holds an ever-growing chunk of the cubic crystal, and the convergence is fast and measured, the effective dimension of a properly extracted cusp slice reaching a clean three at a side of only about ten cells and the gravitational falloff settling toward one over distance not far above it. Infinity is a convenience for the perfect crystal, never a requirement for the physics.

One caution matters for anyone building or picturing the crystal. The rings grown outward by reflection are not themselves horospheres, they are counts of mirror steps from a seed rather than surfaces of equal distance to the ideal point, so growing a few shells and looking will still show a wildly hyperbolic tangle. To see the cubic cusp one must extract a horospherical slice deliberately and read the geometry tangent to it, or build the flat cubic lattice directly as the boundary. Done properly, a finite horosphere chunk already is clean three-dimensional space, while a raw finite reflection patch only approximates it. So the phrase flat space at the edge is exact once read precisely. The flat space is the horosphere, Euclidean at every finite distance, the radial direction across it stays hyperbolic, and the bulk and the cusp are two different places to run the law, the curved interior and the flat boundary, never the same experiment.
